\subsection{Infinitesimal transformations.}

The conception of an ``infinitely small'' transformation goes back to the beginnings of the infinitesimal Calculus; it is known that Descartes discovers the instantaneous centre of rotation by assuming that ``in the infinitely small'' every plane movement can be assimilated into a rotation; the elaboration of analytic Mechanics, in the XVIIIth century, is entirely founded on similar ideas. In 1851, Sylvester, seeking to form invariants of the linear group

GL(3, C) or of certain of its subgroups, gives to the parameters \(z_{j}\) occurring in these matrices ``infinitely small'' increments of the form \(\alpha_{j}dt\) , and expresses that a function \(f((z_{j}))\) is invariant by writing down the equation \(f((z_{j} + \alpha_{j}dt)) = f((z_{j}))\) ; this gives him a linear equation for f with partial derivatives Xf = 0, where

\[
X f = \sum_ {j} \alpha_ {j} \frac {\partial f}{\partial z _ {j}},\tag{25.1}
\]

X being thus a differential operator, ``a derivative in the direction of directed parameters \(\alpha_{j}\) '' ({[}304{]}, vol.~3, pp.~326 and 327); Sylvester seems to feel that there is there a general principle of fairly wide scope, but does not seem to have returned to the question. A little later, Cayley ({[}58{]}, v. II, pp.~164-178) proceeds in the same way for invariants of SL(2, C) in certain representations of this group and shows that it is the solution of two partial differential equations of the first order Xf = 0, Yf = 0, where X and Y are obtained as above starting from ``infinitely small'' transformations

\[
\left( \begin{array}{c c} 0 & 0 \\ d t & 0 \end{array} \right) \text {   and   } \left( \begin{array}{c c} 0 & d t \\ 0 & 0 \end{array} \right).
\]

In modern terms, this is explained by the fact that X and Y generate the Lie algebra \(\mathfrak{sl}(2,\mathbf{C})\) ; besides Cayley calculates explicitly the bracket XY-YX and shows that it also comes from an ``infinitely small'' transformation.

In his memoir of 1868 on groups of movements {[}174 b{]}, Jordan uses from beginning to end the concept of ``infinitely small transformation'', but exclusively from a geometric point of view. There is no doubt that it is with him that the idea of a group with one parameter ``generated'' by an infinitely small transformation appears: for Jordan, it is the set of transformations obtained by ``repeating suitably'' the infinitely small transformation (loc. cit. p.~243). Klein and Lie, in their memoir of 1871, use the same expression ``repeated infinitely small transformation'' ({[}182{]}, v. I, pp.~424-459), but the context shows that they mean by that an integral of a differential system. If the group with one parameter that they consider is made up of transformations \(x' = f(x, y, t)\) , \(y' = g(x, y, t)\) , the corresponding ``infinitely small transformation'' is given by

\[
d x = p (x, y) d t, d y = q (x, y) d t
\]

where \(p(x,y)=\frac{\partial f}{\partial t}(x,y,t_{0}),q(x,y)=\frac{\partial g}{\partial t}(x,y,t_{0})\) , where \(t_{0}\) corresponds to the identical transformation of the group. As Klein and Lie know the functions f and g explicitly, they have no problem in checking that the functions

\[
t \mapsto f (x, y, t) \text {   and   } t \mapsto g (x, y, t)
\]

give in a parametric form the integral curve of the differential equation

\[
q (\xi , \eta) d \xi = p (\xi , \eta) d \eta
\]

going through the point \((x,y)\) , but do not give any general reason for that; besides they do not use this fact again in the rest of their memoir.

